\lysh{34317}{4}

\begin{alist}
\item 
\begin{rlist}
\item Αν κάποιος λείπει από το σπίτι του έχει καταναλώσει $x = 0$ κυβικά μέτρα νερού.
Επειδή $0 \in [0, 30]$ θα χρησιμοποιήσουμε τον τύπο $f(x) = 12 + 0,5x$, οπότε το ποσό
που θα πληρώσει είναι
$f(0) = 12 + 0,5 \cdot 0 = 12$ ευρώ.

\item Επειδή $10 \in [0, 30]$ θα χρησιμοποιήσουμε τον τύπο $f(x) = 12 + 0,5x$, οπότε το ποσό
που θα πληρώσει είναι
$f(10) = 12 + 0,5 \cdot 10 = 12 + 5 = 17$ ευρώ.

\item Επειδή $50 \in (30, +\infty)$ θα χρησιμοποιήσουμε τον τύπο $f(x) = 0,7x + 6$, οπότε το
ποσό που θα πληρώσει είναι
$f(50) = 0,7 \cdot 50 + 6 = 35 + 6 = 41$ ευρώ.
\end{rlist}

\item Διακρίνουμε δύο περιπτώσεις:
\begin{itemize}
\item Ο κάτοικος της πόλης Α κατανάλωσε $x$ κυβικά μέτρα νερού με $x \in [0, 30]$. Επειδή ο
κάτοικος της πόλης Α πλήρωσε μεγαλύτερο λογαριασμό, θα ισχύει:
\begin{align*}
f(x) > g(x) \Leftrightarrow 12 + 0,5x > 12 + 0,6x \ &\Leftrightarrow \\
\Leftrightarrow 12 - 12 > 0,6x - 0,5x \ &\Leftrightarrow \\
\Leftrightarrow 0 &> 0,1x \Leftrightarrow x < 0,
\end{align*}
Άτοπο διότι $x \in [0, 30]$.
Επομένως ο κάτοικος της πόλης Α δεν μπορεί να κατανάλωσε από 0 έως 30 κυβικά
μέτρα νερού.

\item Ο κάτοικος της πόλης Α κατανάλωσε $x$ κυβικά μέτρα νερού, με $x \in (30, +\infty)$. Επειδή
ο κάτοικος της πόλης Α πλήρωσε μεγαλύτερο λογαριασμό, θα ισχύει:
\begin{align*}
f(x) > g(x) \Leftrightarrow 0,7x + 6 > 12 + 0,6x \ &\Leftrightarrow \\
\Leftrightarrow 0,7x - 0,6x > 12 - 6 \ &\Leftrightarrow \\
\Leftrightarrow 0,1x &> 6 \Leftrightarrow x > 60.
\end{align*}

Επομένως, τόσο ο κάτοικος της πόλης Α όσο και αυτός της πόλης Β κατανάλωσαν
περισσότερα από 60 κυβικά μέτρα νερού.
\end{itemize}
\end{alist}

